\documentclass[book.tex]{subfiles}

\begin{document}

\pagestyle{empty}
\vspace*{36 pt}
\begin{center}\Huge\bfseries\itshape
Part 4: Modern Physics
\end{center}
%\newpage

In Part 1 of this book, we covered basic mathematical foundations. In Part 2, we covered physical fundations. In Part 3, we've introduced you to a variety of techniques for physical machine learning and discussed how to use these techniques for modeling physical systems of interest.

However, the systems we've discussed thus far would be of more interest to the applied scientist than to the practising physicist since they in effect offer improved approximation techniques to aid the numerical analysis of physical systems. Such techniques are of course tremendously useful, especially to practicing applied scientists, but arguably don't provide a new way to look at the foundations of physics.

In this Part, we will argue that the techniques of differentiable programming offer new ways to think about physics itself. We will do so by venturing close to the frontier of modern physics, startin from quantum field theory then moving on to study complex systems such as semiconductors, superconductors, fusion. We will conclude this Part, and the book itself, by venturing to the frontier of modern physics by studying the standard model and then quantum gravity.



\subfile{chapters/part5/superconductors}
%\subfile{chapters/part5/Plasma_Physics}
%\subfile{chapters/part5/general_relativity}

\subfile{chapters/part5/the_standard_model}
\subfile{chapters/part5/String_Theory_And_Quantum_Gravity}

\end{document}